% Standalone manuscript: upload this complete file as main.tex to Overleaf.
% Bibliography is embedded; no external figures, .bib files, or custom class.
\documentclass[11pt,a4paper]{article}
\usepackage[T1]{fontenc}
\usepackage[utf8]{inputenc}
\usepackage{lmodern}
\usepackage[margin=24mm]{geometry}
\usepackage{amsmath,amssymb,amsfonts}
\usepackage{array,booktabs}
\usepackage{microtype}
\usepackage{xurl}
\usepackage[hidelinks]{hyperref}
\setlength{\emergencystretch}{3em}
\setlength{\parskip}{0.2em}
\setlength{\parindent}{1.2em}
\setlength{\abovedisplayskip}{7pt plus 2pt minus 2pt}
\setlength{\belowdisplayskip}{7pt plus 2pt minus 2pt}
\renewcommand{\arraystretch}{1.08}
\newcommand{\R}{\mathbb R}
\newcommand{\indic}{\mathbb I}
\newcommand{\clip}{\operatorname{clip}}
\title{Two-View Image-Based Anthropometry:\\
Profile-Guided Reconstruction and Geometric Measurement}
\author{Author details to be completed}
\date{}
\begin{document}
\maketitle

\begin{abstract}
This paper describes an anthropometry pipeline that converts frontal and profile photographs into metric measurements using supplied stature. Pretrained segmentation and pose estimation provide image evidence for an explicit geometric reconstruction. A profile-guided carve rejects displaced profile runs and limits depth through frontal column thickness. A curvature-dependent waist correction modifies local occupancy, while an elliptical shoulder model approximates an across-back path. The implementation specifies coordinate calibration, reconstruction branches, section conditioning, conic recovery, and local sensitivity calculations. Historical pilot records from three adult male participants contain ten tape comparisons, with a mean absolute error of $2.69$\,cm. Verification identifies a sign discrepancy in conic-axis recovery; the reported pilot values precede its correction. The proposed contribution concerns the coupling of profile filtering with columnwise depth constraints and localized shape correction. Its effect on measurement accuracy remains to be established through corrected, independently collected experiments.
\end{abstract}
\noindent\textbf{Keywords:} image-based anthropometry; visual hull; profile-guided reconstruction; geometric measurement.

\section{Introduction and Related Work}
Photographic anthropometry requires correspondence between image-derived geometry and a specified physical measurement. ISO~8559-1 distinguishes circumferences, breadths, and surface paths used in clothing measurement~\cite{iso8559}. Meunier and Yin evaluated a two-dimensional photographic measurement system~\cite{meunier2000}; Hung \emph{et al.} estimated circumferences from image-derived widths and depths~\cite{hung2004}. Automatic measurement from frontal and lateral silhouettes was subsequently developed by Aslam \emph{et al.}~\cite{aslam2017}. These studies establish substantial prior work for two-view measurement and elliptical section models.

The visual hull represents the volume consistent with observed silhouettes~\cite{laurentini1994visual}. Human-shape methods also include the statistical SMPL representation~\cite{loper2015smpl}, volumetric inference in BodyNet~\cite{varol2018bodynet}, and learned reconstruction from orthographic silhouettes~\cite{ji2019}. The present pipeline uses pretrained perception followed by geometric rules with explicit physical units. Its principal proposed adaptation couples a filtered profile interval to frontal run topology and a columnwise thickness bound. This construction addresses depth replication when a profile silhouette contains arms or displaced limbs. A spatially tapered waist correction provides a further local shape constraint. The shoulder model specializes an established ellipse-perimeter approximation to a defined surface path. These adaptations require ablation to establish an accuracy benefit; the supporting segmentation, fitting, and meshing algorithms retain their established attribution.

The following formulation describes the implementation underlying the historical pilot records. Chest is measured from the profile-guided branch, whereas waist and hip use a corrected conventional hull. A separately rounded profile-guided volume supplies the display mesh. This distinction is necessary because each operation imposes a different section model. Geometric processing uses fixed rules after pretrained perception.

\section{Image Preparation and Metric Alignment}
\subsection{Acquisition, masks, and landmarks}
Inputs are frontal and approximately orthogonal profile RGB images, together with stature $H$ in centimetres. The full crown and feet must be visible; close-fitting clothing and consistent posture reduce silhouette ambiguity. Let $X$ denote lateral position, $Y$ distance downwards from the crown, and $Z$ sagittal position, all in millimetres. A pinhole camera gives $u=f_xx_c/z_c+c_u$ and $v=f_yy_c/z_c+c_v$. Replacing depth by a representative $z_0$ yields the weak-perspective approximation used below. A depth displacement $\Delta z$ changes magnification by approximately $-\Delta z/z_0$. The implementation estimates neither camera intrinsics nor residual yaw, so stature calibration alone does not eliminate perspective error.

The implementation uses MediaPipe~\cite{lugaresi2019mediapipe} and its Pose Landmarker segmentation output~\cite{mediapipeguide}. The heavy pretrained model operates in single-image, single-person mode with detection and presence thresholds of $0.5$. Contiguous RGB arrays are padded on the right with replicated columns to a width divisible by four; the segmentation map is cropped back afterwards. This alignment step preserves image scale. A mismatched map is resized bilinearly. Foreground is defined by the strict probability threshold $p>0.5$, followed by selection of the largest eight-connected component. Enclosed background holes smaller than $2\%$ of foreground area are filled; border-connected background is retained. Missing pose or segmentation output, or foreground occupying less than $2\%$ of the frame, causes rejection. Default mask dilation is zero.

An optional GrabCut baseline uses the foreground/background colour-mixture and boundary formulation of Rother \emph{et al.}~\cite{rother2004grabcut}. Landmark-based initialization constructs a torso polygon, limb capsules, and a head region. With estimated pixel stature $H_p=(y_{\rm ankle}-y_{\rm shoulder})/0.779$, seed erosion is $0.012H_p$ and the probable-foreground and background dilations are $0.045H_p$ and $0.13H_p$. A border ring spans $2\%$ of the smaller image dimension. Five iterations with random seed zero precede a $7\times7$ elliptical closing and the same component cleanup. This baseline is separate from the default learned mask.

Landmarks are refined by a second inference on the silhouette bounding box expanded by $15\%$. Crops whose longer dimension is below 1024 pixels are enlarged bicubically; larger crops retain their resolution. For resize factor $\gamma$, resized crop dimensions $(W_r,H_r)$, and crop origin $(x_0,y_0)$, a normalized landmark $(\xi,\eta)$ maps to $(x_0+W_r\xi/\gamma,\,y_0+H_r\eta/\gamma)$. Shoulder, hip, and ankle pairs (indices 11--12, 23--24, 27--28) determine alignment. Elbows and wrists provide subsequent overlap checks. The fit is rejected when more than three requested landmarks have visibility below $0.3$.

\subsection{Weighted alignment and calibration}
For image points $p_i$, canonical targets $q_i$ place each shoulder, hip, and ankle pair horizontally while preserving its Euclidean pair distance and mean vertical position. Shoulder and hip midpoints define a common spine coordinate; ankle stance and image ordering are preserved. Visibility weights are clipped to $[0.05,1]$ and normalized to sum to one. The similarity fit extends orthogonal Procrustes alignment~\cite{schonemann1966generalized}:
\begin{equation}
 \min_{\sigma>0,R\in SO(2),t}\sum_iw_i\|\sigma Rp_i+t-q_i\|^2.
 \label{eq:procrustes}
\end{equation}
Writing $\bar p=\sum_iw_ip_i$, $\bar q=\sum_iw_iq_i$, and
$K=\sum_iw_i(p_i-\bar p)(q_i-\bar q)^T=U\Sigma V^T$, the solution is
\begin{equation}
 \begin{aligned}
 S&=\operatorname{diag}(1,\det(VU^T)),& R&=VSU^T,\\
 \sigma&=\frac{\operatorname{tr}(\Sigma S)}{\sum_iw_i\|p_i-\bar p\|^2},&
 t&=\bar q-\sigma R\bar p.
 \end{aligned}
\end{equation}
Eliminating translation centres the objective; maximizing its cross term gives $R$, and differentiation with respect to scale gives $\sigma$. Near-zero source spread uses unit scale. Rotation beyond $35^\circ$ or scale outside $[0.6,1.6]$ causes rejection. Masks are warped on the original canvas with bilinear interpolation and rethresholded at $0.5$; landmarks receive the same transform. This global alignment can leave articulation differences and can clip foreground at the canvas boundary.

Calibration follows warping. In each view $v\in\{f,s\}$, let $t_v,b_v$ be the first and last occupied rows and $L_v=b_v-t_v+1$. Set $H_m=10H$ and $s_v=H_m/L_v$ millimetres per pixel. The horizontal origin $c_v$ is the equally weighted mean of nonempty-row silhouette midpoints. Projection is
\begin{equation}
 (u_f,v_f)=(c_f+X/s_f,\,t_f+Y/s_f),\qquad
 (u_s,v_s)=(c_s+Z/s_s,\,t_s+Y/s_s).
 \label{eq:projection}
\end{equation}
Nearest-pixel mask lookup returns zero outside the image. Occupancy arrays use $(X,Y,Z)$ ordering. OpenCV section coordinates are $(\text{column},\text{row})=(Z\text{-index},X\text{-index})$ and must be reversed when recovering metric $(X,Z)$ points. Supplied stature is a common scale input to both views. Continuous geometric measurements therefore satisfy $\Delta C/C\simeq\Delta H/H$; fixed metric thresholds and voxel discretization can change the discrete response.

\section{Reconstruction and Local Shape Constraints}
\subsection{Conventional occupancy and branch selection}
For binary masks $M_f,M_s$, the conventional orthographic occupancy is
\begin{equation}
 O_0(X,Y,Z)=M_f(u_f,v_f)M_s(u_s,v_s).
 \label{eq:hull}
\end{equation}
At a given height, every frontal run pairs with every profile run. With $n$ frontal and $m$ profile intervals, this produces up to $nm$ rectangular regions, including limb-depth replication. Reconstruction first uses a 10\,mm grid over silhouette extrema with 30\,mm lateral and sagittal margins and $0\leq Y<H_m$. The occupied coarse bounding box is expanded by one coarse cell and resampled uniformly at $h=2$\,mm. Features absent from the coarse sample can consequently be missed. The fine stage stores a dense array over this box.

The largest 26-connected volume component is retained. Within successive $XZ$ slices, eight-connected components are tracked greedily by nearest centroid, with a maximum linking distance of 60\,mm. Tracks spanning at least five slices have their centroid coordinates and four half-extents filtered by a cubic Savitzky--Golay fit~\cite{savitzky1964smoothing}, using at most 15 slices. Negative half-extents are clipped to zero. Each original component is intersected with its smoothed bounding rectangle, and the largest volume component is selected again. This operation trims occupancy; it cannot recover missing material.

The resulting conventional hull receives the waist correction below and supplies waist and hip sections. Chest uses an independently constructed, unrounded profile-guided volume, followed by its own section conditioning. A rounded copy of the profile-guided volume is used for display. Failure to construct a usable profile core sends chest and display processing to the corrected conventional hull. Shoulder and hip-to-ankle outputs use masks and landmarks directly.

\subsection{Profile filtering and columnwise depth}
At profile row $y$, select the widest occupied run $[\ell_y,u_y]$, with width $d_y=u_y-\ell_y$ and centre $c_y=(\ell_y+u_y)/2$. Reject missing runs and rows satisfying
\begin{equation}
 d_y>1.8\operatorname{median}(d)\quad\text{or}\quad
 |c_y-\operatorname{median}(c)|>
 \max\{6\operatorname{MAD}(c),40\text{ pixels}\}.
\end{equation}
Here MAD is the unscaled median absolute deviation. A 51-row vertical dilation expands the rejected-row mask by approximately 25 rows in either direction. Accepted endpoints are linearly interpolated, with endpoint clamping, then filtered by a 101-row median and a cubic Savitzky--Golay window of at most 61 odd-numbered rows. Windows shrink or are bypassed for short sequences. No accepted rows causes branch failure. For offsets $r$ in a filtering window, the local polynomial coefficients solve
\begin{equation}
 \hat a=(V^TV)^{-1}V^Tf,\qquad V_{rk}=r^k,\quad k=0,\ldots,3;
 \qquad \hat f(0)=\hat a_0.
\end{equation}
Window lengths are in pixels, so their physical bandwidth depends on $s_s$. Interpolated heights are flagged because rejected observations have been replaced by a smooth depth prior.

Converting filtered endpoints to metric coordinates gives $\zeta_-(Y)$ and $\zeta_+(Y)$. For each fixed frontal column $X$, identify its contiguous vertical foreground runs. If the run containing $Y$ spans $n_Y(X,Y)$ fine-grid samples, define
\begin{equation}
 \begin{aligned}
 T(X,Y)&=h\,n_Y(X,Y),&
 D_p(Y)&=\zeta_+(Y)-\zeta_-(Y),\\
 z_c(Y)&=\tfrac12[\zeta_+(Y)+\zeta_-(Y)],&
 z_\pm(X,Y)&=z_c(Y)\pm\tfrac12\min\{D_p(Y),T(X,Y)\},\\
 O_p(X,Y,Z)&=M_f(u_f,v_f)\indic\{z_-\leq Z\leq z_+\}.&&
 \end{aligned}
 \label{eq:profile}
\end{equation}
The profile supplies a shared sagittal midpoint and available depth. The frontal vertical run supplies a local thickness bound: an approximately horizontal limb has a short vertical extent at its outer columns, even when it joins the torso horizontally. This coupling is the main proposed reconstruction adaptation. Its thickness and midpoint assumptions can misrepresent sloping limbs or displaced anatomy. Endpoint interpolation may also extend occupancy beyond the observed profile mask; $O_p$ is therefore a model-conditioned reconstruction. Computation proceeds in chunks of 64 frontal columns while retaining the dense output volume.

\subsection{Measurement levels and curvature-dependent waist correction}
The frontal torso trace starts from the widest horizontal run and subsequently selects the run nearest the preceding midpoint. Separate left/right arm-separation tests use five-row margins. Missing trace values are interpolated and smoothed by a cubic 21-row fit. With mean shoulder and hip levels $Y_s,Y_h$, the waist proxy is $Y_w=Y_s+0.688(Y_h-Y_s)$, snapped to the nearest available trace row. The factor approximates $(0.818-0.620)/(0.818-0.530)$ from engineering body-height fractions. Chest is evaluated at $Y_c=Y_s+0.35(Y_w-Y_s)$; hip uses the landmark level. These prescribed levels do not search for maximum chest or hip girth.

Each torso half-width is clipped using its 15th percentile in a $\pm40$\,mm neighbourhood; nonpositive widths are rejected. Girth extraction expands these bounds by a factor of $1.15$, whereas shoulder extraction uses the unexpanded trace. Conventional-hull sections also receive a centred depth cap $D_c=\min(D,0.75W)$ when a profile elbow or wrist lies within 90\,mm of the measurement level, or when those landmarks are unavailable. Removal exceeding $1\%$ of original depth generates an advisory. Profile-guided chest receives only lateral torso clipping.

Waist refinement locates a local minimum of each half-width within $\pm15$ trace samples of $Y_w$, then walks along rising widths to obtain an enclosing interval $[Y_0,Y_1]$. Fit $w(Y)\simeq\alpha(Y-Y_w)^2+\beta(Y-Y_w)+\gamma$. The graph curvature is $2\alpha/(1+\beta^2)^{3/2}$; the implementation uses the small-slope magnitude $\kappa=|2\alpha|$. Writing $L=Y_1-Y_0$, the inward displacement is
\begin{equation}
 \begin{aligned}
 R_c&=\kappa^{-1},\qquad
 \theta=2\arcsin\!\left(\clip\!\left(\frac{L}{2R_c},-0.999,0.999\right)\right),\\
 d&=\min\!\left\{\frac{R_c}{2}[1-\cos(\theta/2)],\,0.40w_{\rm loc}\right\}.
 \end{aligned}
 \label{eq:waist}
\end{equation}
Here $w_{\rm loc}$ is the smaller of the endpoint and local waist half-widths. The correction is skipped for prominence below 2\,mm, nonpositive $L$, or $\kappa<10^{-6}$\,mm$^{-1}$. The displacement uses half the circular sagitta; before clipping, its small-curvature limit is $d\simeq\kappa L^2/16$. This attenuated geometric prior is derived from the frontal width trace.

Within $Y_w\pm15$\,mm, sampled Hann weights $\omega_j=\tfrac12[1-\cos(2\pi j/(n+1))]$, $j=1,\ldots,n$, suppress abrupt vertical changes. Sagittal tapering uses $t(z)=\cos[\clip(\pi z/(2z_h),-\pi/2,\pi/2)]$, where $z_h$ is the maximum occupied $|Z|$, bounded below by 1\,mm. At each $(Y,Z)$, the left and right torso limits move inward by $\omega_jd_Lt(z)$ and $\omega_jd_Rt(z)$. Only occupancy inside the traced torso interval expanded by $1.2$ is eligible for deletion; exterior material is retained. The curvature-to-displacement mapping and spatial taper constitute the proposed waist adaptation. Actual section-area change is $\Delta A=h^2(N_{\rm before}-N_{\rm after})$; auxiliary curvature-area expressions used in sensitivity calculations do not replace this voxel count.

\section{Section Measurement and Surface Export}
\subsection{Conditioning and conic recovery}
The nearest $XZ$ slice is selected for each measurement level, giving a level quantization error no greater than $h/2$ within the sampled range. After torso clipping and any conditional depth cap, an OpenCV elliptical closing with nominal radius 4\,mm fills small gaps. The largest eight-connected section component is retained and opened with radius $r=0.9\min(W/2,D/2)$, rounded to grid units with a minimum of two pixels. In set notation, closing is $(S\oplus B)\ominus B$ and opening is $(S\ominus B)\oplus B$. The latter imposes substantial section rounding, which influences the reported circumference.

All contour points are retained, and the contour of largest area is converted to metric $(X,Z)$ coordinates. A Graham scan~\cite{graham1972efficient} deduplicates points, sorts by polar angle, retains the farther collinear point, and removes non-left turns. Complexity is $O(n\log n)$. At least five hull vertices are required. Convexification discards section concavities before fitting.

Let $p_i$ be hull vertices and $u_i=p_i-\bar p=(x_i,z_i)^T$ their mean-centred coordinates. The implemented fixed-constant algebraic fit solves
\begin{equation}
 \begin{gathered}
 Ax^2+Bxz+Cz^2+Dx+Ez+1=0,\qquad
 M_i=(x_i^2,x_iz_i,z_i^2,x_i,z_i),\\
 (M^TM)\phi=-M^T\mathbf1,\qquad \phi=(A,B,C,D,E)^T.
 \end{gathered}
 \label{eq:fit}
\end{equation}
Hull vertices have equal weights. This unrestricted conic fit can return a non-ellipse; ellipse-specific constrained fitting offers an established alternative for comparison~\cite{fitzgibbon1999direct}. With
\begin{equation}
 Q=\begin{pmatrix}A&B/2\\B/2&C\end{pmatrix},\quad l=\binom DE,
 \quad \mu=-\tfrac12Q^{-1}l,\quad
 F=\tfrac14l^TQ^{-1}l-1,
 \label{eq:recover}
\end{equation}
completing the square gives $(u-\mu)^TQ(u-\mu)=F$. Eigenvectors of $Q$ determine orientation; positive ratios $F/\lambda_i$ give squared semi-axes, sorted as $a\geq b>0$. The centre is $\bar p+\mu$. Singular systems, near-degenerate discriminants relative to $A^2+B^2+C^2$, zero eigenvalues, and invalid axis ratios cause rejection. Algebraic residual RMS is dimensionless, despite the historical output key \texttt{rms\_residual\_mm}. The fitted ellipse may extend beyond the processed section.

Ramanujan's approximation~\cite{ramanujan1914modular} supplies the perimeter in millimetres,
\begin{equation}
 P(a,b)=\pi[3(a+b)-\sqrt{(3a+b)(a+3b)}],\qquad C=P(a,b)/10\text{ cm}.
 \label{eq:perimeter}
\end{equation}
The pilot code instead used $F_{\rm impl}=-1-(AE^2+CD^2-BDE)/(4AC-B^2)$ in axis recovery. Equation~\eqref{eq:recover} requires a plus sign before this fraction. The discrepancy vanishes when the fitted linear terms vanish, but mean-centring nonuniform hull samples does not guarantee that condition. A verification case with true axes $(160,100)$\,mm returned $(155.359,97.099)$\,mm under the historical expression; completing the square recovered the true axes. Pilot values below are preserved as historical outputs and do not validate the corrected formulation.

\subsection{Shoulder path and linear proxy}
At the engineering acromion level $Y_a=0.182H_m$, let $W,D$ be unexpanded traced width and profile depth. When $W,D>0$ and $D/W\leq0.6$, shoulder length is $S=\tfrac12P(W/2,D/2)$. This half-perimeter assumes that shoulder endpoints coincide with the transverse extrema of a symmetric section and that the posterior path is elliptical. Its contribution is a specified path model with explicit applicability checks; a tape trajectory can differ from these assumptions.

When this path cannot be computed, frontal silhouette tips are sought between $Y_s-20$\,mm and $Y_s+0.4(Y_e-Y_s)$, where $Y_e$ is the lower available elbow level. Missing elbows give an upper limit $Y_s+0.18(Y_h-Y_s)$. The trace retains the arm-fused prefix, applies a quadratic Savitzky--Golay window of at most nine rows, and selects each half-width maximum. A post-peak taper of at least 5\,mm on each side marks reliable tips. Their two-dimensional Euclidean separation gives projected breadth; a tip-to-landmark breadth ratio outside $[1,1.5]$ is advisory. Otherwise the fallback is the shoulder-landmark chord multiplied by $s_f$. The legacy fallback label can cover either breadth method, so the method must accompany the value. The output historically named inseam is $Y_{\rm ankle}-Y_{\rm hip}$, a vertical hip-to-ankle proxy without a corresponding tape reference in this pilot.

\subsection{Display geometry and retained outputs}
Display rounding is applied to a separate volume. Each $XZ$ component is opened with a Euclidean distance-transform element of radius $0.9$ times its smaller half-extent, with a two-pixel minimum; an empty opening retains the component. Limb seeds satisfy frontal vertical thickness divided by profile-core depth $<0.5$, grow through ratios $<0.85$, and require frontal area at least $2000$\,mm$^2$. Limb $YZ$ openings are united with the rounded $XZ$ result. Distance-transform erosion uses $d_{\rm background}>r$ and dilation uses distance to the eroded set $\leq r$, with padding exceeding $r$. These elements differ from the rasterized OpenCV elements used for section measurement.

Junction repair attempts chessboard-distance closings of 20, 40, 60, and 100\,mm, intersecting the result with the unrounded profile volume. Processing stops at one component or residual components totalling at most $0.5\%$. The largest component is exported, with an advisory if more than $1\%$ is removed. Zero padding of $\lceil3\sigma_g\rceil+1$ voxels precedes Gaussian smoothing with $\sigma_g=1$ voxel and marching cubes at isovalue $0.5$ with spacing $h$~\cite{lorensen1987marching}. Origin and padding offsets are undone, and $(x,y,z)$ maps to $(x,z,y_{\max}-y)$ for a floor-based, vertical-$Z$ display frame. The largest mesh component is retained; a watertight mesh with negative signed volume is inverted. Twelve Taubin passes~\cite{taubin1995signal} use $\lambda=0.5$ and $\nu=-0.53$ under the applicable library convention. Optional simplification targets 150\,000 faces.

Saved outputs include normalized masks, conventional and display occupancy, and measurement records containing levels, fitted axes, rotations, and shoulder method. Advisory messages are largely logged and are not all serialized into the historical JSON. Reproducibility therefore requires retaining the diagnostic log with image/model identifiers and numerical settings. Dense Boolean occupancy consumes approximately $N_XN_YN_Z$ bytes; each full double-precision coordinate array adds eight times this storage. Chunked evaluation limits temporary arrays while leaving dense output storage unchanged.

\section{Sensitivity and Interpretation of Error}
First-order propagation gives $\operatorname{Var}(C)\simeq J\Sigma J^T$ for differentiable inputs with covariance $\Sigma$~\cite{jcgm2008}. The implemented reporting index combines scale, boundary, curvature, conic-fit, and depth-cap terms as
\begin{equation}
 \delta_C^2=\delta_s^2+\delta_b^2+\delta_c^2+\delta_e^2+\delta_d^2.
 \label{eq:rss}
\end{equation}
This sum neglects covariance among terms derived from the same masks. Computation uses millimetres and converts the final index to centimetres. It is a local sensitivity indicator with incomplete model-error coverage.

With assumed localization spread $\sigma_p=1.5$ pixels and inclusive spans $L_v$, set $r_v=\sqrt2\sigma_p/L_v$. The implementation uses $\delta_s=P\sqrt{r_f^2+r_s^2}$, with perimeter $P$ in millimetres. For $q=(3a+b)(a+3b)$,
\begin{equation}
 P_a=\pi\!\left(3-\frac{3a+5b}{\sqrt q}\right),\qquad
 P_b=\pi\!\left(3-\frac{5a+3b}{\sqrt q}\right),\qquad
 \delta_b=\frac{\sigma_p(s_f+s_s)}2\sqrt{P_a^2+P_b^2}.
 \label{eq:derivative}
\end{equation}
An axis-specific scale model would propagate $(P_aa r_f)^2+(P_bb r_s)^2$ under independent view-scale assumptions. Uncertainty in supplied stature would add a common contribution $(P\sigma_H/H)^2$, omitted from the pilot index. Neither expression accounts for camera perspective or garment thickness. Under a regular convex parallel-boundary offset $\tau$, the perimeter change is $2\pi\tau$; its use as a pixel-error model requires a corresponding physical offset assumption. A one-pixel offset at 2\,mm/pixel would then change perimeter by approximately 12.57\,mm.

For conic residuals $r=M\phi+\mathbf1$ and more than five vertices, the algebraic fit covariance is approximated by
\begin{equation}
 \hat\sigma^2=\frac{r^Tr}{n-5},\qquad
 \Sigma_\phi=\hat\sigma^2(M^TM)^{-1},\qquad
 \delta_e=\sqrt{\max(J_\phi\Sigma_\phi J_\phi^T,0)}.
\end{equation}
The perimeter Jacobian uses central differences with parameter step $\max(10^{-4}|\phi_i|,10^{-10})$. Five-point fits or unavailable covariance return zero in the historical helper, as do derivatives whose perturbed conics fail. These cases represent unavailable sensitivity information and should be distinguished from negligible uncertainty in subsequent implementations. Algebraic residual covariance also depends on the fixed conic normalization and correlated hull sampling.

The waist helper uses the auxiliary quantity $A_{\rm miss}=dL-\tfrac12d^2\tan(\theta/2)$ and curvature spread $\sigma_\kappa=2\sqrt{\operatorname{Var}(\alpha)}$. Its perimeter contribution is
\begin{equation}
 \delta_c=\sqrt{\frac{\pi}{A_{\rm corr}}}
 \left[\sum_{j\in\{L,R\}}
 \left(\frac{\partial A_{{\rm miss},j}}{\partial\kappa_j}\sigma_{\kappa,j}\right)^2\right]^{1/2},
 \label{eq:curvunc}
\end{equation}
where $A_{\rm corr}$ is corrected section area. Differentiation holds $L$ fixed and uses curvature step $\max(10^{-3}|\kappa|,10^{-9})$. Because $L$ is a vertical trace span, $A_{\rm miss}$ is an auxiliary model quantity rather than the removed $XZ$ area. The historical helper also centres its analysis on the narrowest trace row, which can differ from the proportional waist level used for correction. These differences restrict its interpretation as a physical error estimate.

For depth removed by conditional capping, the allowance is $\delta_d=|P_b|\Delta D/2$. Interpolated profile depth has no separate uncertainty term. Linear lengths use $\delta_L^2=(Lr_f)^2+2(\sigma_ps_f)^2$. This frontal-only expression was also assigned to the two-view shoulder output. For the shoulder path, a branch-consistent approximation would instead use $\operatorname{Var}(S)\simeq(P_a\sigma_W/4)^2+(P_b\sigma_D/4)^2$, augmented by covariance and path-model error. Level placement adds $\Delta C\simeq C'(Y)\Delta Y$ where the section varies smoothly. Discrete component changes, clipping, and branch failure require empirical assessment beyond local derivatives.

\section{Historical Pilot Results and Limitations}
The pilot comprises three adult male project authors, of stature 172, 177, and 173\,cm. Their image dimensions were $1086\times1448$, $960\times1280$, and $900\times1600$ pixels, respectively. Acquisition was unblinded, without independently repeated capture, repeated operators, or camera calibration. Table~\ref{tab:pilot} reproduces the available comparisons from the historical implementation. Ranges recorded for participant S3's hip and shoulder were represented by their midpoints. No tape reference was available for the hip-to-ankle proxy.

\begin{table}[htbp]
\centering
\caption{Historical pilot comparisons, in centimetres. Error is estimate minus reference. The index $\delta$ follows Eq.~\eqref{eq:rss}; it is not a calibrated confidence interval. Estimates precede the conic-recovery correction.}
\label{tab:pilot}
\small
\begin{tabular}{llrrrr}
\toprule
Participant & Measurement & Reference & Estimate & $\delta$ & Error\\
\midrule
S1 & Chest    & 92.50 & 87.52 & 1.01 & $-4.98$\\
S1 & Waist    & 85.00 & 82.25 & 1.01 & $-2.75$\\
S1 & Hip      & 98.00 &100.61 & 1.03 & $+2.61$\\
S2 & Chest    & 89.00 & 90.36 & 1.35 & $+1.36$\\
S2 & Waist    & 94.00 & 93.01 & 1.35 & $-0.99$\\
S2 & Shoulder & 48.00 & 44.71 & 0.43 & $-3.29$\\
S3 & Chest    & 93.98 & 88.15 & 1.11 & $-5.83$\\
S3 & Waist    & 85.09 & 82.05 & 1.10 & $-3.04$\\
S3 & Hip      & 98.43 & 99.32 & 1.12 & $+0.89$\\
S3 & Shoulder & 47.62 & 48.81 & 0.37 & $+1.19$\\
\bottomrule
\end{tabular}
\end{table}

Across these ten comparisons, mean absolute error is 2.69\,cm, root-mean-square error 3.13\,cm, signed bias $-1.48$\,cm, and mean absolute percentage error $3.39\%$. The maximum absolute error is 5.83\,cm. The eight girth comparisons alone have mean absolute error 2.81\,cm. Repeated body sites within participants are correlated; pooled summaries therefore describe these records and cannot establish population performance. Only two references fall within estimate $\pm\delta$, consistent with incomplete error coverage. Three segmentation repetitions for each of the six images were pixel-identical in the tested environment. This demonstrates same-input repeatability under those conditions; acquisition repeatability remains unmeasured.

The shoulder-path estimates were 44.71 and 48.81\,cm against references of 48.00 and 47.62\,cm, while landmark chords were 36.52 and 30.39\,cm. These examples illustrate the consequence of using a surface-path model and do not establish its general validity. Agreement studies require consistent measurement definitions and sufficiently independent observations~\cite{bland1986}. Corrected end-to-end results are needed before attributing discrepancies to the geometric priors, especially because conic-axis recovery affects both estimates and conic-derived sensitivity.

The main unresolved question is whether the profile/thickness coupling improves measurement at occluded levels across postures and body shapes. An ablation should compare conventional occupancy and filtered profile occupancy with and without the column bound, while holding perception, level selection, and section fitting fixed. A corresponding waist comparison should separate curvature displacement from section morphology, since both alter perimeter. The half-ellipse shoulder path requires tape endpoints matched to its assumed transverse extrema. Independent participants and repeated acquisition are necessary to separate systematic geometric error from mask and capture variation. The present pilot cannot support a claim of general accuracy or incremental improvement for any proposed component.

\section{Conclusion}
The pipeline specifies metric anthropometry from two silhouettes through profile-core reconstruction constrained by frontal column thickness and a localized curvature-dependent waist adjustment. Branch-specific section extraction and a shoulder-path model expose the geometric assumptions. Historical pilot results require conic-recovery correction and independent validation before they can establish accuracy or the value of the proposed adaptations.

\begin{thebibliography}{99}
\small
\setlength{\itemsep}{0.1em}
\setlength{\parskip}{0pt}
\bibitem{iso8559}
International Organization for Standardization, \emph{Size designation of clothes---Part 1: Anthropometric definitions for body measurement}, ISO 8559-1:2017, 2017. [Online]. Available: \url{https://www.iso.org/standard/61686.html}.

\bibitem{meunier2000}
P.~Meunier and S.~Yin, ``Performance of a 2D image-based anthropometric measurement and clothing sizing system,'' \emph{Applied Ergonomics}, vol.~31, no.~5, pp.~445--451, 2000. doi: \href{https://doi.org/10.1016/S0003-6870(00)00023-5}{10.1016/S0003-6870(00)00023-5}.

\bibitem{hung2004}
P.~C.-Y.~Hung, C.~P.~Witana, and R.~S.~Goonetilleke, ``Anthropometric measurements from photographic images,'' in \emph{Work with Computing Systems 2004}, H.~M.~Khalid, M.~G.~Helander, and A.~W.~Yeo, Eds. Kuala Lumpur: Damai Sciences, 2004, pp.~764--769. [Online]. Available: \url{https://www.ieda.ust.hk/dfaculty/ravi/papers/wwcs2004_1.pdf}.

\bibitem{aslam2017}
M.~Aslam, F.~Rajbdad, S.~Khattak, and S.~Azmat, ``Automatic measurement of anthropometric dimensions using frontal and lateral silhouettes,'' \emph{IET Computer Vision}, vol.~11, no.~6, pp.~434--447, 2017. doi: \href{https://doi.org/10.1049/iet-cvi.2016.0406}{10.1049/iet-cvi.2016.0406}.

\bibitem{laurentini1994visual}
A.~Laurentini, ``The visual hull concept for silhouette-based image understanding,'' \emph{IEEE Transactions on Pattern Analysis and Machine Intelligence}, vol.~16, no.~2, pp.~150--162, 1994.

\bibitem{loper2015smpl}
M.~Loper, N.~Mahmood, J.~Romero, G.~Pons-Moll, and M.~J.~Black, ``SMPL: A skinned multi-person linear model,'' \emph{ACM Transactions on Graphics}, vol.~34, no.~6, article~248, 2015.

\bibitem{varol2018bodynet}
G.~Varol, D.~Ceylan, B.~Russell, J.~Yang, E.~Yumer, I.~Laptev, and C.~Schmid, ``BodyNet: Volumetric inference of 3D human body shapes,'' in \emph{Proceedings of ECCV}, pp.~20--36, 2018.

\bibitem{ji2019}
Z.~Ji, X.~Qi, Y.~Wang, G.~Xu, P.~Du, X.~Wu, and Q.~Wu, ``Human body shape reconstruction from binary silhouette images,'' \emph{Computer Aided Geometric Design}, vol.~71, pp.~231--243, 2019. doi: \href{https://doi.org/10.1016/j.cagd.2019.04.019}{10.1016/j.cagd.2019.04.019}.

\bibitem{lugaresi2019mediapipe}
C.~Lugaresi et al., ``MediaPipe: A framework for building perception pipelines,'' arXiv:1906.08172, 2019.

\bibitem{mediapipeguide}
Google AI Edge, ``Pose landmark detection guide.'' [Online]. Available: \url{https://ai.google.dev/edge/mediapipe/solutions/vision/pose_landmarker}. Accessed October 17, 2026.

\bibitem{rother2004grabcut}
C.~Rother, V.~Kolmogorov, and A.~Blake, ``GrabCut: Interactive foreground extraction using iterated graph cuts,'' \emph{ACM Transactions on Graphics}, vol.~23, no.~3, pp.~309--314, 2004.

\bibitem{schonemann1966generalized}
P.~H.~Sch\"onemann, ``A generalized solution of the orthogonal Procrustes problem,'' \emph{Psychometrika}, vol.~31, pp.~1--10, 1966.

\bibitem{savitzky1964smoothing}
A.~Savitzky and M.~J.~E.~Golay, ``Smoothing and differentiation of data by simplified least squares procedures,'' \emph{Analytical Chemistry}, vol.~36, no.~8, pp.~1627--1639, 1964.

\bibitem{graham1972efficient}
R.~L.~Graham, ``An efficient algorithm for determining the convex hull of a finite planar set,'' \emph{Information Processing Letters}, vol.~1, no.~4, pp.~132--133, 1972.

\bibitem{fitzgibbon1999direct}
A.~Fitzgibbon, M.~Pilu, and R.~B.~Fisher, ``Direct least square fitting of ellipses,'' \emph{IEEE Transactions on Pattern Analysis and Machine Intelligence}, vol.~21, no.~5, pp.~476--480, 1999.

\bibitem{ramanujan1914modular}
S.~Ramanujan, ``Modular equations and approximations to $\pi$,'' \emph{Quarterly Journal of Mathematics}, vol.~45, pp.~350--372, 1914.

\bibitem{lorensen1987marching}
W.~E.~Lorensen and H.~E.~Cline, ``Marching cubes: A high resolution 3D surface construction algorithm,'' \emph{ACM SIGGRAPH Computer Graphics}, vol.~21, no.~4, pp.~163--169, 1987.

\bibitem{taubin1995signal}
G.~Taubin, ``A signal processing approach to fair surface design,'' in \emph{Proceedings of SIGGRAPH}, pp.~351--358, 1995.

\bibitem{jcgm2008}
Joint Committee for Guides in Metrology, \emph{Evaluation of measurement data---Guide to the expression of uncertainty in measurement}, JCGM~100:2008. [Online]. Available: \url{https://doi.org/10.59161/JCGM100-2008E}.

\bibitem{bland1986}
J.~M.~Bland and D.~G.~Altman, ``Statistical methods for assessing agreement between two methods of clinical measurement,'' \emph{The Lancet}, vol.~327, no.~8476, pp.~307--310, 1986. doi: \href{https://doi.org/10.1016/S0140-6736(86)90837-8}{10.1016/S0140-6736(86)90837-8}.
\end{thebibliography}
\end{document}
